\newpage
\section{Fazit und Ausblick}

\subsection{Fazit}

Der im Prompt bereitgestellte \ac{UI}-Kontext beeinflusst die Qualität der generierten Tests für WebGIS-Anwendungen deutlich und wirkt sich dabei auf zwei getrennte Aspekte aus. Die Ausführbarkeit der Tests hängt vor allem davon ab, ob das Modell die Elemente der Oberfläche überhaupt findet. Dieses Problem wird größtenteils schon durch den Accessibility-Snapshot behoben. Ob die generierten Tests aber auch inhaltlich das Richtige prüfen, hängt dagegen mit dem Zugriff auf das Kartenmodell zusammen. Erst mit den Map-Model-Helfern lässt sich der Kartenzustand verlässlich abfragen und erst dann steigt die inhaltliche Qualität der Karteninteraktionsprüfungen an (vgl. Kapitel~6.3). Die Dimensionen \emph{coverage} und \emph{assertion} bleiben davon dagegen fast unberührt und über alle Stufen nahezu konstant. Sie hängen also stärker vom Verständnis des Use Cases ab als vom bereitgestellten \ac{UI}-Kontext (vgl. Kapitel~6.3). Der eigentliche Hebel ist also nicht einfach nur \enquote{mehr Kontext}, sondern dem \ac{LLM} gezielt die Mittel bereitzustellen, die es für das Auffinden der Elemente und die Prüfung des Kartenzustands benötigt.

Zwischen der automatisch generierten und manuell erstellten \ac{UI}-Map liegt kaum ein Unterschied in der Testqualität, dafür aber ein erheblicher Unterschied in Erstellungsaufwand (vgl. Kapitel~6.1, 7.4). Der zusätzliche Pflegeaufwand einer manuell erstellten \ac{UI}-Map lohnt sich für diese Anwendung damit nicht.

An seiner Grenze ist der bereitgestellte Kontext aber noch nicht angekommen. Was auch der Stufe~4 fehlt, ist Wissen über Lage, Größe und Überlagerung von Elementen. Die Karte lässt sich zwar über das Kartenmodell inhaltlich abfragen, bleibt aber ein Canvas-Element ohne direkte \ac{DOM}-Struktur. Bedienpanels, die über der Kartenfläche liegen, führen deshalb wiederholt zu Fehlklicks, weil sie die Stellen auf der Karte verdecken, auf die geklickt werden soll. Auch die Umrechnung einer Kartenkoordinate in eine Klickposition gelingt ohne passende Hilfsfunktion nur unzuverlässig.

Auch die iterative Rückmeldung aus der Testausführung stößt an eine Grenze. Sie hebt die Ausführbarkeit noch einmal deutlich an, allerdings auf Kosten von Aufwand und Ressourcen (vgl. Kapitel~6.4) und bringt eine Schwäche mit. Ein Loop, der erst bei einem bestandenen Ergebnis abbricht, optimiert auf dieses Kriterium und nicht auf die Korrektheit der Prüfung. In einem produktiven Umfeld wäre dieses Verhalten sehr riskant (vgl. Kapitel~7.1).

Der ergänzende Lauf mit einem zweiten Modell bestätigt den beobachteten Kontexteffekt. Das stützt die Annahme, dass es sich tatsächlich um einen Effekt des Kontexts handelt und nicht um eine Eigenheit eines einzelnen Modells.

Ob ein generierter Test besteht, sagt wenig darüber aus, ob er auch das prüft, was der Use Case fordert. Das zeigt sich besonders deutlich in den unteren Kontextstufen, in denen ein großer Anteil bestandener Tests auf einer unzureichenden Prüfung beruht (vgl. Kapitel~6.3). Erst die Kombination aus deterministischer Ausführung und semantischer Bewertung macht diesen Unterschied sichtbar. Eine vollständig automatische Generierung, ohne dass jemand nachprüft, was ein Test tatsächlich tut, bleibt damit auch mit umfangreichem \ac{UI}-Kontext nicht zuverlässig möglich.

\subsection{Ausblick}

Die Ergebnisse eröffnen mehrere Anschlusspunkte für zukünftige Arbeiten. Dazu gehört zunächst die Erweiterung der Untersuchungsumgebung. Die Demo-Anwendung ist als Single-Page-Anwendung umgesetzt. Reale WebGIS-Portale, wie etwa Klima- oder Umweltportale, verteilen ihre Funktionen häufig über mehrere Seiten. Seitenübergreifende und deutlich anspruchsvollere Use Cases stellen den Workflow nochmal vor zusätzliche Herausforderungen. Zudem bleiben Karteninteraktionen außerhalb des sichtbaren Kartenausschnitts unberücksichtigt. Wenn ein Feature erst durch Verschieben oder Zoomen der Karte zugänglich wird, kommt neben der Koordinatenumrechnung noch eine Navigationsplanung hinzu, die der Kontext bislang nicht unterstützt. Offen ist auch wie sich die Befunde auf andere WebGIS-Frameworks wie Leaflet oder MapLibre übertragen lassen.

Auch in der Methodik ist noch Spielraum für Anpassungen am Generierungs-Workflow. Der hier untersuchte Loop verwendet ein einzelnes Modell, das Generierung und Korrektur in einem Schritt übernimmt. Playwright stellt mit seinen Test Agents mittlerweile eine rollenbasierte Alternative bereit: Ein Planner erkundet die Anwendung und erstellt einen Testplan, ein Generator setzt diesen Plan in Testcode um, und ein Healer prüft fehlschlagende Tests zur Laufzeit und korrigiert sie \parencite{microsoft_agents_nodate}. Diese Agenten erkunden die Anwendung dabei selbst, statt auf einen vorab erfassten Kontext angewiesen zu sein, und reihen sich damit in den von \textcite[S.~81]{kim_towards_2026} beschriebenen Ansatz ein, bei dem Agenten Testaufgaben selbstständig planen und ausführen, statt starren Skripten zufolgen. Ob eine solche Aufteilung in getrennte Rollen die beobachteten Schwächen beheben würde und wie sie mit GIS-spezifischen Anforderungen umgehen, wäre ein naheliegender Anschlusspunkt.

Der letzte Punkt betrifft die Erfassung von Aufwand und Kosten. In dieser Arbeit wurde nur die Gesamtlaufzeit je Stufe erfasst, nicht aber Token- oder Zeitverbrauch jedes einzelnen Generierungsaufrufs. Gerade für die Bonus-Stufe~5 mit der hohen Anzahl an Aufrufen wäre interessant, an welcher Stelle im Loop der größte Aufwand anfällt.
